\chapter{DISCUSSION}

\section{Principal findings}

This thesis found a clear mismatch between HPV vaccine awareness and self-reported any-dose HPV vaccination among Vietnamese women aged 15--29 years. Nearly two-thirds of women had heard of HPV vaccination, but fewer than one in thirteen had been vaccinated. Among women who were already aware, only 12.0\% had received the vaccine. This pattern demonstrates that information alone is not sufficient to generate uptake.

The second major finding is the strong socioeconomic gradient. Awareness was higher among women in richer households, urban residents, Kinh/Hoa women, and women with higher education. Vaccination showed an even steeper relative gradient: 16.6\% of women in the richest quintile reported vaccination compared with only 0.7\% in the poorest quintile. The standard concentration index confirmed that vaccination had stronger relative socioeconomic concentration than awareness. The Erreygers indices should be interpreted more cautiously as absolute inequality measures for bounded outcomes because they are influenced by outcome prevalence.

The third major finding is that the awareness--vaccination gap was concentrated among poorer aware women. This distinction matters. A woman who has never heard of the vaccine requires information and communication. A woman who has heard of the vaccine but remains unvaccinated may face financial, geographic, service delivery, social, or prioritization barriers. The gap outcome therefore points toward access and affordability, not only health education.

\section{Interpretation of associated factors}

Wealth was the most consistent marker across the pathway. In adjusted models, the poorest women had substantially lower awareness and vaccination than the richest women. They were also more likely to remain unvaccinated despite awareness. The decomposition results reinforced this conclusion: the poorest quintile was the largest contributor to both awareness and vaccination inequality. Because MICS6 predates the planned public financing of HPV vaccination, the observed socioeconomic gradient should be interpreted as a baseline measure of inequity under a predominantly opportunistic or privately financed access context. This is consistent with the access framework of Penchansky and Thomas, in which affordability, availability, and accessibility all shape use of health services \cite{penchansky1981}.

Education also played an important role. Women with primary education or less had markedly lower awareness in the adjusted model. Lower secondary education was associated with lower vaccination and a higher awareness--vaccination gap. Education may operate through health literacy, confidence in preventive care, ability to interpret vaccine information, and capacity to navigate services \cite{nutbeam2008}.

Age patterns suggest possible missed catch-up opportunities. Vaccination prevalence was lowest among women aged 15--19 years and highest among those aged 20--24 years, which may reflect delayed vaccination after adolescence, differences in family decision-making, private-market affordability, or contact with reproductive health services. Because MICS6 is cross-sectional and does not record age at vaccination, these mechanisms cannot be separated. Programmatically, however, the pattern supports monitoring not only routine target-age vaccination but also older adolescents and young adults who may be missed during the transition to public financing.

Rural residence and ethnic minority status were strongly associated with lower awareness and vaccination in bivariate analysis. Ethnic minority status should be interpreted as a marker of structural marginalization, geographic access, language barriers, service availability, and historical inequities rather than as an intrinsic cause of lower vaccination. After adjustment, ethnic minority status remained associated with lower awareness but not vaccination, suggesting that some of the vaccination difference may be mediated through wealth, education, residence, and other structural variables. Nevertheless, ethnic minority women had very low vaccination prevalence in descriptive analysis, and this remains programmatically important.

Media and ICT exposure showed different roles. ICT exposure was almost universal and therefore not suitable for adjusted regression. Mass media exposure was strongly associated with awareness in bivariate analysis, but its adjusted association was weaker. In the decomposition, media exposure contributed meaningfully to inequality, suggesting that unequal distribution of information channels still matters even when other factors are controlled.

The gender-equitable attitude measure was based on rejecting all five wife-beating justifications. It was associated with higher awareness after adjustment. This variable should not be interpreted as direct decision-making autonomy; rather, it may capture broader gender norms, empowerment, and openness to preventive health information.

\section{Comparison with previous evidence}

The findings align with international evidence that HPV vaccination coverage and any-dose uptake are shaped by financing, delivery strategy, and social determinants \cite{bruni2021,gallagher2018}. Countries with organized and publicly financed programs have achieved larger and more equitable population impact, while reliance on individual payment and opportunistic uptake can produce inequitable coverage. WHO's 2024 product-choice guidance and single-dose updates are relevant for Viet Nam because simplified schedules can reduce procurement, service delivery, and follow-up barriers that often disadvantage poorer and remote communities \cite{who2024productchoice,who2024singledose}. The results also support previous methodological work showing that odds ratios can exaggerate associations for common outcomes; therefore, prevalence ratios are preferable for cross-sectional analyses of common outcomes such as awareness \cite{barros2003,zou2004}.

The magnitude of the awareness--vaccination gap in this thesis highlights a specific implementation problem. Raising awareness remains necessary, especially among poor, rural, less educated, and ethnic minority women. However, awareness campaigns alone will not close the vaccination gap unless they are combined with affordable and accessible vaccination services.

\section{Policy implications}

The results support an equity-oriented HPV vaccination strategy with four components. Resolution 104/NQ-CP and the 2026--2028 EPI implementation plan place HPV vaccination on Viet Nam's pathway for planned inclusion and public financing in the Expanded Programme on Immunization from 2026 \cite{govvn2022nq104,baochinhphu2025epi}. The implication is not simply that public financing is needed, but that implementation should be designed so that public financing does not reproduce the socioeconomic gradient observed in the MICS6 baseline.

First, financing matters. Public financing or strong subsidy mechanisms are likely to be necessary if Viet Nam aims to prevent HPV vaccination from functioning primarily as a private-market service accessible to better-off households.

Second, delivery should prioritize geographic access. Rural and remote areas require delivery channels that reduce travel time and indirect costs. School-based, community-based, or primary health care-based delivery can be considered depending on local implementation capacity.

Third, communication should be targeted rather than only universal. Women in poorer households, ethnic minority communities, and lower education groups had lower awareness. Communication materials should be culturally appropriate, understandable, and available through trusted community and health system channels.

Fourth, monitoring should include equity indicators. National or provincial reporting should not only track overall doses or any-dose uptake, but also uptake by wealth, ethnicity, education, residence, province/region, and school attendance. Equity monitoring is essential for cervical cancer elimination goals \cite{who2020strategy}. Particular attention is needed for poor households, rural communes, ethnic minority communities, out-of-school girls, and young women who are older than the initial school-based target age.

\section{Strengths and limitations}

This study used nationally representative MICS6 data and accounted for the complex survey design in all main analyses. The analysis also moved beyond simple prevalence estimation by using concentration indices and decomposition to quantify socioeconomic inequality.

Several limitations should be considered. First, the cross-sectional design does not permit causal inference; terms such as associated factors, correlates, and components of inequality pattern in this thesis refer to statistical and decomposition patterns, not proven causal effects. Second, HPV vaccination status was self-reported and may be affected by recall or reporting error; it measures receipt of any HPV vaccine dose, not completed vaccination schedule or full protection. Third, because detailed information on dose number, age at vaccination, vaccine type, price paid, provider type, and reasons for non-vaccination was unavailable, the analysis could not distinguish incomplete vaccination, delayed vaccination, affordability barriers, service availability constraints, and vaccine hesitancy. Although cost and service availability were not directly measured, the observed wealth and residence gradients are consistent with access-related barriers. Fourth, the full-sample vaccination outcome coded women who had never heard of HPV vaccination as unvaccinated, consistent with the skip pattern; this is appropriate for estimating reported any-dose vaccination in the full population under a conservative assumption, but it combines awareness and uptake. The vaccination-among-aware analysis was therefore added to separate uptake conditional on awareness. Fifth, the age range 15--29 years reflects older adolescents and young adult women rather than the usual 9--14-year primary target group, so the estimates should not be interpreted as target-age girl coverage. Sixth, sexual experience was self-reported in a face-to-face survey and may be affected by underreporting, particularly among unmarried young women in settings where premarital sex is socially sensitive. Seventh, ICT exposure was nearly universal and could not be reliably included in adjusted models. Eighth, the household wealth index is an asset-based relative measure and may not fully capture disposable income, vaccine affordability, or urban--rural differences in cost of access. Ninth, decomposition results depend on model specification; because no uncertainty intervals were estimated for individual contributions and wealth quintiles were included in a wealth-ranked decomposition, rankings should be interpreted as approximate descriptive accounting rather than causal attribution or statistically tested differences between contributors.

\clearpage
